\chapter{33}

\section{Samstag, 25. August}



Im Haus war es still.

Louis sank zurück in die Kissen.

Der Traum wog schwer. Es dauerte, bis das Zittern nachliess, bis sein
Atem ruhiger wurde und sich sein Körper entspannte. Kaum fand
er zu einer Spur Leichtigkeit zurück, rammte ihn die Nacht wie einen
Pflock zurück in den Schlamm. Selbst im Schlaf verfolgte Giorgia ihn. 

Zerknirscht schlief er knapp vor Sonnenaufgang wieder ein.

\sectionbreak

Es war nach zehn, als ihn das Licht weckte. Der Tag schien jede Ecke mit Leben zu füllen. Louis blinzelte sich von
der Decke den Wänden entlang in den Raum. Von unten vernahm er plötzlich
lauten, rockigen Sound. Steve Lees Stimme erkannte er an den ersten
Tönen. Mein Gott, \emph{«Child in Time»}, durchfuhr es ihn, was für
eine Überraschung. Dieser Song hier oben? Das
Gitarrensolo kannte er beinahe auswendig. Während er sich auszog und die
Tür zur kleinen Eckdusche des \emph{Schäferlofts} öffnete, war Lees
Stimme kaum mehr hörbar. Sie wurde offensichtlich vom Geräusch eines
Staubsaugers übertönt. Das Holzhaus war auffallend schalldurchlässig.
Hier blieb wohl nichts verborgen. Der Gedanke liess ihn schaudern.

\qrblock{Jon Lord \& Steve Lee \emph{«Child in Time»}}{media/QR-Codes/029_Child_in_Time_-_Live_2009_Stay_Tuned,_Jon_Lord_Spotify.png}

Als Louis die Treppe hinuntergestiegen in der Küche stand, fühlte er
sich erstaunlich ausgeruht. Trotz des nachwirkenden Alptraums. Sein
körperlicher Zustand erinnerte ihn an eine passable Version seiner
selbst.

Im hinteren Teil der Küche schob Joh die Staubsaugerbürste mit kräftigem
Griff über ein Liegesofa. Das Möbel nahm mit seiner beeindruckenden Grösse
locker einen Viertel des Raumes ein. Die quadratische Form war speziell.
Es sah erlesen und modern aus. Auch hier im
Wohnzimmerbereich zeigte sich die Einrichtung als harmonische
Kombination von alt und neu. Louis war entzückt. 
Er blieb im Küchenbereich stehen und sah Joh zu. Selbst in Schlabberhose und mit
zerzaustem Haar ging eine faszinierende Eleganz von ihm aus.
Joh musste Louis im Augenwinkel gesehen haben. Er drehte sich um und zog
blitzschnell den Kabelstecker.

«Ciro, hi, guten Morgen auf der \emph{Sunnewaid},» er breitete die Arme
aus, «Kaffee, Hunger?»

Louis nickte begeistert. Schon stand Joh hinter der Kaffeemaschine, die unter kräftigem Brummen eine Portion Pulver in
den Kolben mahlte. Louis setzte sich in Johs Nähe an das Ende des langen, massiven Holztisches.

«Schön habt ihr’s hier oben,» begann er, als Joh ihm eine grossbauchige
Tasse Schaumkaffee hinstellte, «und mein \emph{Schäferloft} ist der
Knaller, ich habe gut geschlafen.»

Joh hob den Daumen, strahlte Louis an und stellte dann eine Köstlichkeit
nach der andern auf den Tisch. Der reichhaltige Brunch erschien ihm einmal mehr wie eine Wiedergutmachung nach
Tagen willkürlichen Futterns.

«Schlag zu, Ciro.»

Joh biss in ein Stück Käse. Für einen guten Moment assen sie stumm. Und
Louis kam vor, als habe Joh ewig Zeit. Nichts anderes zu tun, als hier
mit ihm zu sitzen und zu frühstücken.

«Hab ich dich mit der Musik geweckt?»

Zwischen zwei Kaffeeschlucken liess Louis die Schultern fallen.

«Ne, ne, ich war wach, doch mit Steve Lee geweckt zu werden, wäre nicht
verkehrt gewesen, unvergleichlich, diese Stimme.»

«Da hast du recht. Ich weiss noch genau, wie uns sein tödlicher Unfall
traf, Manuela und mich. Er und \emph{Gotthard} haben für uns noch heute
eine grosse Bedeutung,» Johs Augen blitzten, «wir haben an einem
Gotthard-Konzert unser Einjähriges gefeiert. Ist ein paar Jährchen her.»

Joh legte eine Atempause ein. Er sah aus, als schwelge er für einen
Moment in alten Zeiten.
Louis rechnete nach.

«Dann seid ihr schon viele Jahre zusammen, du und Manuela?»

«Ja, wir haben es tatsächlich schon auf über zehn Jahre Miteinander
geschafft,» Joh setzte sich gerade auf und sah stolz aus, «grossartige
und auch genauso turbulente Jahre. Ohne Manuela, ich weiss nicht, wo ich
gerade herumirren würde.»

Jetzt blickte Joh in die Ferne. Er schien ebenso redselig wie Manuela.
Louis war je länger desto verblüffter. Wie kamen die beiden nur dazu, ihm so viel Persönliches zu erzählen.

«Sie fand das zwar albern, sie sprach lieber von Fügung, was ich noch heute nicht ganz verstehe.»

Er rührte in der Tasse.

«Fügung?»

Louis runzelte die Stirn.

«Hm, dass Menschen sich begegnen, die was miteinander zu tun haben, und
dass es ein Glück ist, wenn sie einander erkennen und mutig genug sind, sich aufeinander einzulassen, 
so in der Art.»

Joh schob seine Brille zurecht und blickte Louis verschwörerisch an.

«Na ja, Frauen ticken anders als wir Männer. Gerade darum gefallen sie
mir. Wobei, bei euch jungen Leuten mag sich das verändert haben, ich
weiss nicht, wie siehst du das, Ciro?»

Louis fühlte sich überfahren. Am liebsten hätte er Joh einfach weiter
zugehört. Er sammelte sich.

«Das hab ich mir noch nie so genau überlegt, weiss nicht, ich kenn die
Welt deiner Generation zu wenig. Klar, ich hab meine Eltern erlebt, aber
als Kind habe ich mir solche Fragen nicht gestellt. Da war alles so, wie
es war und eigentlich ziemlich kompliziert.»

Louis spürte, wie sich sein Hals zusammenzog. Mehr wollte er ihm nicht
verraten. Joh nickte kaum sichtbar. Dann klopfte er Louis sanft auf die
Schulter und stand auf.

«Genug gefrühstückt?» fragte er mit heller Stimme und rückte seinen 
Stuhl nach hinten.

«Ja, ich bin satt bis oben, das war herrlich.»

«Dann führ ich dich doch mal durch Haus und Scheune, ja?»

Louis folgte Joh die Treppe hinauf ins Obergeschoss. Beim Hochsteigen
fiel ihm ein, wie galant Joh Manuela behandelte, sie gestern geradezu
hofiert hatte. Die beiden hatten den Eindruck gemacht, ein frisch
verliebtes Paar zu sein. Und Joh sprach von einer zehnjährigen Beziehung? 
Das war aussergewöhnlich oder zumindest interessant. Sollte
er hier wirklich bleiben können, würde er Joh später bei Gelegenheit
fragen, wie sie das schafften. Giorgia und er hatten es nicht geschafft.
Der Gedanke drückte ihm aufs Herz.
Rechts des Flurs lagen drei, links vier Räume. Joh öffnete als Erstes
die breiteste Tür rechts.

«Hier \emph{Die Oase}, unser Wellnessraum, wenn du so willst. Wenn du
dir mal Zeit für dich nehmen möchtest.»

Vor Louis tat sich ein übergrosses Badezimmer auf. In der Mitte erhöht
und über zwei Treppenstufen erreichbar, eine Wanne, wie er keine
vergleichbare je gesehen hatte. Blass hellblau und glänzend, mit einigen
erhöhten Sitzflächen an den Rändern und überall komplett abgerundet, lag
sie einladend da. Seine Sprachlosigkeit schien Joh zu gefallen. Louis ging näher heran und entdeckte eingebaute Wasserdüsen. Die Wanne bot Platz für mindestens vier Personen.

«Der Wahnsinn,» entfuhr es Louis, mehr als überrascht, «wo gibt es denn
so was?»

Joh schmunzelte. Sein Blick verriet, dass er Reaktionen wie diese genoss.

«Die haben wir selbst gemacht, gemauert, ausgekleidet und gespritzt,
Manuela und ich. Von der Idee bis zur Ausführung haben wir Blut
geschwitzt. Wir haben Tage damit verbracht, die Rundungen zu schleifen,
mein Gott, wenn wir gewusst hätten, was uns erwartet, ich weiss nicht - »
Joh rollte die Augen, «aber es hat sich gelohnt. Wir sind sehr zufrieden
mit unserm Werk.»

Dann führte er Louis weiter durch den Flur. Er erklärte ihm die Nutzung
der Räume. Ans Bad grenzte Manuelas Zimmer, mit \emph{Lady M.}
angeschrieben. Dann folgte das \emph{Schäferloft.} Die Türe gegenüber
war die nächste und letzte, die Joh öffnete. Louis sah darin Wände mit
Gestellen und einer Menge Wäsche. Geradeaus an der Schmalseite des
Raumes stand ein Möbelstück, abgedeckt mit einem losen, weissen Tuch.
Die Form deutete auf ein abgestelltes Klavier hin.

«Hier, Ciro, ist unser Wäscheraum, und hier - ,» Joh zeigte auf die
Beschriftung oberhalb der nächsten Tür, « - hier ist unser \emph{Aletsch,} das
ist aktuell Mazis Zimmer.»

Mit einem Mal blieb Joh stehen. Er hob seinen Arm und stützte sich mit der Hand an der Wand ab. Er grinste.

«Ich rede schon, als hätten wir uns auf deine Anstellung geeinigt, etwas
voreilig, verzeih, aber weisst du, wie du einfach so auftauchst, nach
unserer verzweifelten Suche nach Verstärkung, ist fast ein Wunder.»

Wenn der wüsste. Joh hatte seine Rettung gerade vollkommen in der Hand.

«Ach, schon okay, nach euren Erzählungen kann ich dich gut verstehen.
Und vielleicht finden wir ja zu einer Zusammenarbeit.»

Joh klopfte Louis auf die Schulter.

«Uff, gut, dann also weiter. \emph{Mister Joh} bin dann ich. Die
Beschriftungen waren übrigens Manuelas Vorschlag,» fuhr er fort, schob
wiederholt seine Brille zurecht, und zeigte auf die letzte Tür, «und
hier im \emph{Boxenstopp} sind unsere Jungs einquartiert. Der Raum ist
grösser, hat Platz für zwei, ähnlich wie das \emph{Schäferloft}.»

«Danke für die Führung, das Haus ist geräumiger, als ich von aussen
dachte und sehr aussergewöhnlich eingerichtet, Kompliment, richtig cool
habt ihr es hier.»

Joh winkte Louis hinter sich her.
Die angebaute Scheune erreichten sie über eine direkte Tür neben der
Treppe. An einen grossen Vorraum, der der Unterbringung verschiedener
Geräte diente und Arbeitskleidung an den Wänden hängen hatte, grenzte 
eine nächste Tür. Sie führte in einen riesigen Stall.

«Hier wohnen unsere Schafe, wenn sie von der Alp runter kommen.»

Joh streckte seinen Arm aus und zeigte von links nach rechts.

«Ein bisschen verwaist gerade, unser Stall. Manchmal kann ich mir im
Sommer kaum mehr vorstellen, wie lebendig und laut es hier werden kann.»

Die Grösse der Scheune beeindruckte Louis. Und auch, wie ordentlich hier
alles seinen Platz hatte.

«Nun, Ciro, das Wichtigste hast du gesehen, es gibt noch allerhand mehr,
aber das klären wir step by step - doch, warte, einen Raum will ich dir
noch zeigen.»

Joh deutete auf eine weitere Tür. Sie sah robuster aus als alle andern und war abgeschlossen.

«Sollte mal etwas passieren, ist gut, wenn du weisst, dass wir uns
erstmal vor Ort darum kümmern.»

Joh zog einen Schlüsselbund aus der Hosentasche und schloss auf.
Was sich hier wohl verbarg? Joh drückte die Klinke hinunter und machte
das Licht an.
Überraschungen würde Louis hier oben wohl noch einige erleben. Es war
definitiv an der Zeit, seine Vorstellungen über ein \emph{Leben auf dem
Bauernhof} über Bord zu werfen. Auf einen Operationssaal, oder eine
Arztpraxis mit OP-Tisch oder wie sich dieser kleine Saal auch nennen
mochte, war er zuletzt vorbereitet. Der rundum weiss gekachelte Raum mit
einem hellgrau ausgelegten Linoleumboden, mit Arzneischränken an den
Wänden, einem Schreibtisch und zwei Stühlen aus Chrom dahinter, reinlich
sauber alles, kamen Louis, an einem Ort wie diesem, fremd vor. Er
schaute Joh fragend an.
Die Hofpraxis sei ein Puzzle ihres Konzeptes. Johs Erklärungen
hörten sich wie kleine Vorträge an. Louis hing ihm nach wie vor an den Lippen.

«Grundsätzlich sind wir medizinische Selbstversorger. Wenn immer möglich
behandeln wir kleinere, je nach dem auch grössere Unfälle und Krankheiten
unserer Tiere hier vor Ort. Auch unser Team und einige Nachbarn
inklusive deren Tiere werden bei uns ärztlich versorgt. Eine
zusätzliche, kleine Einkommensquelle ist daraus geworden, neben unserer
Schafzucht und meiner Zusammenarbeit mit dem Kanton. Und \emph{den} Teil
hat dir ja Manuela bereits erklärt.»

Joh lehnte sich an die Wand.

«Ach so, du bist auch Arzt?» fragte Louis interessiert.

Joh winkte ab.

«Nein, nein, die Medizinerin bei uns ist Manuela. Ich bin ursprünglich Psychiatriepfleger. 
Das Drumherum stiess mir irgendwann
auf, schlechte Bedingungen und so, ich habe dann eine zweite Ausbildung 
zum Sozialarbeiter gemacht.»

«Aha, Manuela ist Tierärztin?»

«So ungefähr,» Joh holte Luft, «sie hängte an ein Medizinstudium später
eine mehrjährige Weiterbildung am rechtsmedizinischen Institut an,
im Gebiet Rechts- und Gerichtsmedizin. Ihr Ziel als junge Frau war die
forensische Pathologie,» Joh schmunzelte, «sie wollte für Gerechtigkeit kämpfen,» jetzt malte er Anführungsstriche in die Luft, «ja, wenn Manuela mal brennt, dann aber richtig,» seine Augen funkelten, «ich hatte grosses Glück, sie damals kennen zu lernen.»

Louis war beeindruckt.

«Du machst mich neugierig. Und \emph{wie} hast du sie denn
kennengelernt?»
